\chapter{Combinatorial}

\section{Permutations}  
	\subsection{Cycles}
		Let $g_S(n)$ be the number of $n$-permutations whose cycle lengths all belong to the set $S$. Then
		$$\sum_{n=0} ^\infty g_S(n) \frac{x^n}{n!} = \exp\left(\sum_{n\in S} \frac{x^n} {n} \right)$$

	\subsection{Derangements}
		Permutations of a set such that none of the elements appear in their original position.
		\[ \mkern-2mu D(n) = (n-1)(D(n-1)+D(n-2)) = n D(n-1)+(-1)^n = \left\lfloor\frac{n!}{e}\right\rceil \]

	\subsection{Burnside's lemma}
		Given a group $G$ of symmetries and a set $X$, the number of elements of $X$ \emph{up to symmetry} equals
		 \[ {\frac {1}{|G|}}\sum _{{g\in G}}|X^{g}|, \]
		 where $X^{g}$ are the elements fixed by $g$ ($g.x = x$).

		 If $f(n)$ counts ``configurations'' (of some sort) of length $n$, we can ignore rotational symmetry using $G = \mathbb Z_n$ to get
		 \[ g(n) = \frac 1 n \sum_{k=0}^{n-1}{f(\text{gcd}(n, k))} = \frac 1 n \sum_{k|n}{f(k)\phi(n/k)}. \]

\section{Partitions and subsets}
	\subsection{Partition function}
		Number of ways of writing $n$ as a sum of positive integers, disregarding the order of the summands.
		\[ p(0) = 1,\ p(n) = \sum_{k \in \mathbb Z \setminus \{0\}}{(-1)^{k+1} p(n - k(3k-1) / 2)} \]
		\[ p(n) \sim 0.145 / n \cdot \exp(2.56 \sqrt{n}) \]

		\begin{center}
		\begin{tabular}{c|c@{\ }c@{\ }c@{\ }c@{\ }c@{\ }c@{\ }c@{\ }c@{\ }c@{\ }c@{\ }c@{\ }c@{\ }c}
			$n$    & 0 & 1 & 2 & 3 & 4 & 5 & 6  & 7  & 8  & 9  & 20  & 50  & 100 \\ \hline
			$p(n)$ & 1 & 1 & 2 & 3 & 5 & 7 & 11 & 15 & 22 & 30 & 627 & $\mathtt{\sim}$2e5 & $\mathtt{\sim}$2e8 \\
		\end{tabular}
		\end{center}

	\subsection{Stars and bars}
		Solutions of $x_1 + \dots + x_k = n$: $\binom{n+k-1}{k-1}$ with $x_i \ge 0$, $\binom{n-1}{k-1}$ with $x_i \ge 1$.
		With $0 \le x_i \le c$: $\sum_{i} (-1)^i \binom{k}{i} \binom{n - i(c+1) + k - 1}{k - 1}$.

	\subsection{Lucas' Theorem}
		Let $n,m$ be non-negative integers and $p$ a prime. Write $n=n_kp^k+...+n_1p+n_0$ and $m=m_kp^k+...+m_1p+m_0$. Then $\binom{n}{m} \equiv \prod_{i=0}^k\binom{n_i}{m_i} \pmod{p}$.

\section{General purpose numbers}
	\subsection{Bernoulli numbers}
		EGF of Bernoulli numbers is $B(t)=\frac{t}{e^t-1}$ (FFT-able), which gives $B_1=-\frac12$ (used below; BernoulliNumber.h uses $+\frac12$).
		$B[0,\ldots] = [1, -\frac{1}{2}, \frac{1}{6}, 0, -\frac{1}{30}, 0, \frac{1}{42}, \ldots]$

		Sums of powers:
		\small
		\[ \sum_{i=1}^n i^m = \frac{1}{m+1} \sum_{k=0}^m \binom{m+1}{k} B_k \cdot (n+1)^{m+1-k} \]
		\normalsize

		Euler-Maclaurin formula for infinite sums:
		\small
		\[ \sum_{i=m}^{\infty} f(i) = \int_m^\infty f(x) dx - \sum_{k=1}^\infty \frac{B_k}{k!}f^{(k-1)}(m) \]
		\[ \approx \int_{m}^\infty f(x)dx + \frac{f(m)}{2} - \frac{f'(m)}{12} + \frac{f'''(m)}{720} + O(f^{(5)}(m)) \]
		\normalsize

	\subsection{Stirling numbers of the first kind}
		Number of permutations on $n$ items with $k$ cycles.
		\begin{align*}
			&c(n,k) = c(n-1,k-1) + (n-1) c(n-1,k),\ c(0,0) = 1 \\
			&\textstyle \sum_{k=0}^n c(n,k)x^k = x(x+1) \dots (x+n-1)
		\end{align*}
		$c(8,k) = 0, 5040, 13068, 13132, 6769, 1960, 322, 28, 1$ \\
		$c(n,2) = 0, 0, 1, 3, 11, 50, 274, 1764, 13068, 109584, \dots$

	\subsection{Eulerian numbers}
		Number of permutations $\pi \in S_n$ in which exactly $k$ elements are greater than the previous element. $k$ $j$:s s.t. $\pi(j)>\pi(j+1)$, $k+1$ $j$:s s.t. $\pi(j)\geq j$, $k$ $j$:s s.t. $\pi(j)>j$.
		$$E(n,k) = (n-k)E(n-1,k-1) + (k+1)E(n-1,k)$$
		$$E(n,0) = E(n,n-1) = 1$$
		$$E(n,k) = \sum_{j=0}^k(-1)^j\binom{n+1}{j}(k+1-j)^n$$

	\subsection{Motzkin numbers}
		\begin{itemize}
		    \item different ways of drawing non-intersecting chords between n points on a circle.
		    \item \textbf{positive} integer sequences of length $n + 1$ in which the opening and ending elements are 1, and the difference between any two consecutive elements is -1, 0 or 1.
		    \item $1, 1, 2, 4, 9, 21, 51, 127, 323, 835, ...$
		\end{itemize}
		\begin{align*}
		    M_n = M_{n-1} + \sum_{i=0}^{n-2} M_i M_{n-2-i} = \frac{2n+1}{n+2} M_{n-1} + \frac{3n-3}{n+2} M_{n-2}
		\end{align*}

	\subsection{Delannoy numbers}
		\begin{itemize}
		    \item Number of ways to move from (0,0) to (m,n) with 3 kinds of steps: (1,0), (1,1), (0,1)
		    \item $D(0, k) = 1, 1, 1, 1, ...$
		    \item $D(3, k) = 1, 7, 25, 63, ...$
		\end{itemize}
		\begin{align*}
		    D(m,n) = \sum_{k=0}^{\min(m,n)} \binom{m+n-k}{m} \binom{m}{k} = \sum_{k=0}^{\min(m,n)} \binom{m}{k} \binom{n}{k} 2^k
		\end{align*}

	\subsection{Stirling numbers of the second kind}
		Partitions of $n$ distinct elements into exactly $k$ groups.
		$$S(n,k) = S(n-1,k-1) + k S(n-1,k)$$
		$$S(n,1) = S(n,n) = 1$$
		$$S(n,k) = \frac{1}{k!}\sum_{j=0}^k (-1)^{k-j}\binom{k}{j}j^n$$

	\subsection{Bell numbers}
		Number of partitions of a set of $n$ items into non-empty disjoint subsets. $B(n) =$
		$1, 1, 2, 5, 15, 52, 203, 877, 4140, 21147, \dots$. For $p$ prime,
		\begin{align*}
		    B_{n+1} = \sum_{k=0}^{n} \binom{n}{k} B_k\\
		    B_{p+n} \equiv B_n + B_{n+1} \pmod{p}\\
		    B_{p^m+n} \equiv mB_n + B_{n+1} \pmod{p}
		\end{align*}

	\subsection{Labeled unrooted trees}
		\# on $n$ vertices: $n^{n-2}$ \\
		\# on $k$ existing trees of size $n_i$: $n_1n_2\cdots n_k n^{k-2}$ \\
		\# with degrees $d_i$: $(n-2)! / ((d_1-1)! \cdots (d_n-1)!)$

	\subsection{Catalan numbers}
		\[ C_n=\frac{1}{n+1}\binom{2n}{n}= \binom{2n}{n}-\binom{2n}{n+1} = \frac{(2n)!}{(n+1)!n!} \]
		\[ C_0=1,\ C_{n+1} = \frac{2(2n+1)}{n+2}C_n,\ C_{n+1}=\sum C_iC_{n-i} \]
		${C_n = 1, 1, 2, 5, 14, 42, 132, 429, 1430, 4862, 16796, 58786, \dots}$
		\begin{itemize}[noitemsep]
			\item sub-diagonal monotone paths in an $n\times n$ grid.
			\item strings with $n$ pairs of parenthesis, correctly nested.
			\item full binary trees with with $n+1$ leaves (0 or 2 children).
			\item ordered trees with $n+1$ vertices.
			\item ways a convex polygon with $n+2$ sides can be cut into triangles by connecting vertices with straight lines.
			\item permutations of $[n]$ with no 3-term increasing subseq.
			\item correct bracket sequence of length $2n$.
		\end{itemize}


\section{Graphs}
	\subsection{Kirchhoff's matrix-tree theorem}
		Laplacian $L = D - A$ (multi-edges count, loops ignored). The number of spanning trees is any cofactor of $L$:
		delete row and column $i$ and take the determinant. Directed, out-arborescences rooted at $r$:
		$L = D_{in} - A$, delete row and column $r$. Weighted edges give the sum over trees of the weight product.
	\subsection{LGV lemma}
		Sources $a_1..a_n$, sinks $b_1..b_n$ in a DAG, $M_{ij}$ = (weighted) number of paths $a_i \to b_j$.
		$\det M$ = signed sum over vertex-disjoint path systems; if the only possible matching is $a_i \to b_i$
		(e.g. planar grid, sorted endpoints) it counts the non-intersecting systems.

\section{Games}
	\subsection{Sprague--Grundy}
		$g(s) = \mathrm{mex}\{g(t) : s \to t\}$, terminal states have $g = 0$. A position loses iff $g = 0$.
		Sum of independent games: $g = g_1 \oplus g_2 \oplus \dots$
	\subsection{Nim variants}
		Nim: first player wins iff $\bigoplus a_i \ne 0$. Misère Nim: if all $a_i \le 1$ first player wins iff
		$\bigoplus a_i = 0$, otherwise as normal Nim. Staircase Nim (move coins down one step, ground = step 0): Nim on steps 1, 3, 5, \dots
		Take at most $k$ per move: $g = a \bmod (k + 1)$. Moore's Nim$_k$ (move in up to $k$ piles): write piles
		in binary; first player loses iff every bit column sum is $\equiv 0 \pmod{k+1}$.

\section{Useful numbers}
	\textbf{Perfect} ($= 2^{p-1}(2^p - 1)$, $2^p - 1$ prime): 6, 28, 496, 8128, 33550336, 8589869056, 137438691328, 2305843008139952128.
	Mersenne prime exponents $p$: 2, 3, 5, 7, 13, 17, 19, 31, 61, 89, 107, 127.

	\textbf{Most divisors} $\max_{n \le N} d(n)$ (attained at): $10^3$: 32 (840), $10^4$: 64 (7560), $10^5$: 128 (83160),
	$10^6$: 240 (720720), $10^7$: 448 (8648640), $10^8$: 768 (73513440), $10^9$: 1344 (735134400),
	$10^{12}$: 6720 (963761198400), $10^{15}$: 26880, $10^{18}$: 103680 (897612484786617600).

	\textbf{Distinct prime factors}: $\le 9$ for $n \le 10^9$ (223092870), $\le 15$ for $n \le 10^{18}$.

	\textbf{$\pi(10^k)$}: 4, 25, 168, 1229, 9592, 78498, 664579, 5761455, 50847534 ($k = 1..9$).
	\textbf{Max prime gap}: 220 below $10^8$, 282 below $10^9$, 1442 below $10^{18}$.

	\textbf{Primes}: $10^9 + 7$, $10^9 + 9$, $10^{18} + 3$, $10^{18} + 9$, $2^{61} - 1$, $2^{63} - 25$ (largest $< 2^{63}$).
	NTT $c \cdot 2^k + 1$ (primitive root): 998244353 $= 119 \cdot 2^{23} + 1$ (3), 167772161 $= 5 \cdot 2^{25} + 1$ (3),
	469762049 $= 7 \cdot 2^{26} + 1$ (3), 754974721 $= 45 \cdot 2^{24} + 1$ (11), 2013265921 $= 15 \cdot 2^{27} + 1$ (31).
	Primitive root of $10^9 + 7$: 5.

	\textbf{Unlabeled trees} ($n = 1..15$), rooted: 1, 1, 2, 4, 9, 20, 48, 115, 286, 719, 1842, 4766, 12486, 32973, 87811;
	free: 1, 1, 1, 2, 3, 6, 11, 23, 47, 106, 235, 551, 1301, 3159, 7741.

	\textbf{Fits in ll}: $20!$, $\binom{66}{33}$, $F_{92}$ (not $21!$, $\binom{67}{33}$, $F_{93}$).

\kactlimport{PartitionNumbers.h}